\documentclass[a4paper]{article}

\usepackage{amssymb}
\usepackage{amsmath}
\usepackage{mathtools}


%opening
\title{Daniel's H2 Mathematics Summaries}
\author{Daniel Ginting}
\date{\today}

\begin{document}
	
	\begin{titlepage}
		\thispagestyle{empty}
		
		\begin{center}
			{\Large Singapore-Cambridge GCE A-Level}\\[0.5cm]
			{\large H2 Mathematics (9758)}\\[3cm]
			
			{\Huge\bfseries H2~Mathematics~Summaries}\\[1cm]
			{\Large 2025/2026}\\[3cm]
			
			{\large Daniel Ginting}\\[0.3cm]
			{\large Yishun Innova Junior College}\\[2cm]
			
			{\large \today}
		\end{center}
		
		\vfill
		
		\begin{center}
			\emph{-- Soli Deo Gloria --}
		\end{center}
		
	\end{titlepage}
	
	
%\maketitle

\newpage
\thispagestyle{empty}

\vspace*{\fill}
\begin{center}
	\emph{To poki.}
\end{center}
\vspace*{\fill}

\newpage

\begin{abstract}

This set of summaries has been devised specifically according to YIJC's lecture notes, cohort 2025/2026,  Singapore A-Levels Syllabus code 9758. This set is by no means a replacement for nor a subsitute to the lecture notes themselves. In fact, the corresponding lecture notes are a prerequisite to understanding and using this document. It serves as a reference to contents and concepts which otherwise have been understood, and therefore is not a learning material by itself.

\end{abstract}


\newpage

% nA
\setcounter{section}{0}
\renewcommand{\thesection}{1A}
\section{Standard Graphs}

Lorem ipsum dolor sit amet, consectetur adipiscing elit. Sed do eiusmod tempor
incididunt ut labore et dolore magna aliqua. Ut enim ad minim veniam, quis
nostrud exercitation ullamco laboris nisi ut aliquip ex ea commodo consequat.
Duis aute irure dolor in reprehenderit in voluptate velit esse cillum dolore eu
fugiat nulla pariatur. Excepteur sint occaecat cupidatat non proident, sunt in
culpa qui officia deserunt mollit anim id est laborum.

% nB
\renewcommand{\thesection}{1B}
\section{Conic Sections}

Lorem ipsum dolor sit amet, consectetur adipiscing elit. Sed do eiusmod tempor
incididunt ut labore et dolore magna aliqua. Ut enim ad minim veniam, quis
nostrud exercitation ullamco laboris nisi ut aliquip ex ea commodo consequat.
Duis aute irure dolor in reprehenderit in voluptate velit esse cillum dolore eu
fugiat nulla pariatur. Excepteur sint occaecat cupidatat non proident, sunt in
culpa qui officia deserunt mollit anim id est laborum.

\renewcommand{\thesection}{1C}
\section{Parametric Curves}

Lorem ipsum dolor sit amet, consectetur adipiscing elit. Sed do eiusmod tempor
incididunt ut labore et dolore magna aliqua. Ut enim ad minim veniam, quis
nostrud exercitation ullamco laboris nisi ut aliquip ex ea commodo consequat.
Duis aute irure dolor in reprehenderit in voluptate velit esse cillum dolore eu
fugiat nulla pariatur. Excepteur sint occaecat cupidatat non proident, sunt in
culpa qui officia deserunt mollit anim id est laborum.

% reset
\renewcommand{\thesection}{\arabic{section}}
\setcounter{section}{1}


% nA
\setcounter{section}{1}
\renewcommand{\thesection}{2A}
\section{Basic Properties of Vectors in Two and Three Dimensions}

Lorem ipsum dolor sit amet, consectetur adipiscing elit. Sed do eiusmod tempor
incididunt ut labore et dolore magna aliqua. Ut enim ad minim veniam, quis
nostrud exercitation ullamco laboris nisi ut aliquip ex ea commodo consequat.
Duis aute irure dolor in reprehenderit in voluptate velit esse cillum dolore eu
fugiat nulla pariatur. Excepteur sint occaecat cupidatat non proident, sunt in
culpa qui officia deserunt mollit anim id est laborum.

% nB
\renewcommand{\thesection}{2B}
\section{Scalar and Vector Products in Vectors}

Lorem ipsum dolor sit amet, consectetur adipiscing elit. Sed do eiusmod tempor
incididunt ut labore et dolore magna aliqua. Ut enim ad minim veniam, quis
nostrud exercitation ullamco laboris nisi ut aliquip ex ea commodo consequat.
Duis aute irure dolor in reprehenderit in voluptate velit esse cillum dolore eu
fugiat nulla pariatur. Excepteur sint occaecat cupidatat non proident, sunt in
culpa qui officia deserunt mollit anim id est laborum.

\renewcommand{\thesection}{2C}
\section{(Vector) Lines}

\renewcommand{\thesection}{2D}
\section{(Vector) Planes}

% reset
\renewcommand{\thesection}{\arabic{section}}
\setcounter{section}{2}

\section{Transformation of Graphs}

Transformations, and the phrasings:
\begin{itemize}
	\item Translation
	
	A translation of $n$ units in the \emph{positive} / \emph{negative} $x$/$y$-direction.
	\item Scaling
	
	A scaling with a scale factor of $n$ parallel to the $x$/$y$-axis.
	\item Reflection
	
	A reflection in the $x$/$y$-axis.
\end{itemize}

Manipulations (remember the technicalities of absolute and reciprocal):
\begin{itemize}
	\item Absolute
	\item Reciprocal
\end{itemize}



\section{Inequalities and Equations}

Odd power factor: alternate

Even power factor: remain

Rational: put all on one side, do not cross multiply.

\[
|x| < a \implies -a < x < a
\]
\[
|x| > a \implies x < -a \lor x > a
\]

\section{Functions}

Notations:

\[
[-\infty, \infty)
\]
\[
\{ x \in \mathbb{R} : a < x < b \}
\]

The range of $f$ is found by graphing.

Horizontal line test phrasing:
\begin{quote}
	"If any horizontal line $y = k$ $(k \in \mathbb{R})$ intersects the graph of
	$y = f(x)$ \emph{at most once}, then $f$ is a one-one (or one-to-one) function."
\end{quote}

$f^{-1}$ exists if $f$ is one-one (or one-to-one).

To find $f^{-1}$, make $x$ the subject.

\[
\text{Domain of } f^{-1} = \text{Range of } f
\]

Conversely,
\[
\text{Range of } f^{-1} = \text{Domain of } f
\]
 
$ y = f^{-1}(x) $ is a reflection of $ y = f(x) $ about $ y = x $.

$fg$ exists if 
\[
\text{Range of } g \subseteq \text{Domain of} f
\]

\[
\text{Domain of } fg = \text{Domain of } g.
\]

Self-inverse function, by Mr Kai, not in the notes:
\begin{itemize}
	\item Same rules and same domain.
	\item Same horizontal and vertical asymptotes.
	\item Symmetry about the line $y = x$.
	\item $f^{2}(x) = x$.
\end{itemize}

\setcounter{section}{5}
\renewcommand{\thesection}{6A}
\section{Introduction to Sequences and Series}

Sequence is list.
Series is sum of list.

Remember how to compute recurrence relation using GC.
\[
u_1 = S_1
\]
\[
u_n = S_n - S_{n-1}, \text{for } n \ge 2
\]

Convergence: 
\[
S_{\infty} = \lim_{n \to \infty} S_n
\]

\renewcommand{\thesection}{6B}
\section{Arithmetic Progression and Geometric Progression (AP\&GP)}

AP:
\[
u_n = a + (n-1)d
\]
\[
d = u_n - u_{n-1}
\]
\[
S_n = \frac{n}{2}[2a + (n-1)d]
\]
\[
u_n - u_{n-1} = c, \text{where } c \text{ is a constant}, \text{ for all } n \in \mathbb{Z}^{+}, n \ge 2
\]

GP:
\[
u_n = ar^{n-1}
\]
\[
r = \frac{u_n}{u_{n-1}}
\]
\[
S_n = \frac{ a ( 1 - r^n ) }{ 1 - r } = \frac{ a( r^n - 1 ) }{ r - 1 }
\]
\[
\lim_{ n \to \infty } S_n = \frac{a}{1-r}
\]
\[
S_\infty \text{ exists } \iff |r| < 1
\]
\[
\frac{u_n}{u_{n-1}} = k, \text{where } k \text{ is a non-zero constant, for all } n \in \mathbb{Z}^{+} , n \ge 2 
\]


\renewcommand{\thesection}{6C}
\section{Sigma Notation}

\renewcommand{\thesection}{\arabic{section}}
\setcounter{section}{6}

\section{Differentiation Techniques}

\section{Applications of Differentiation}

\section{Integration Techniques}

\section{Applications of Integration}

\section{Maclaurin Series}

\section{Differential Equations}

\section{Complex Numbers}

\section{Permutations and Combinations}

\section{Probability}

\setcounter{section}{15}
\renewcommand{\thesection}{16A}
\section{Discrete Random Variable}

\renewcommand{\thesection}{16B}
\section{Binomical Distribution}

\renewcommand{\thesection}{\arabic{section}}
\setcounter{section}{16}

\section{Normal Distribution}

\section{Sampling Distribution}

\section{Hypothesis Testing}

\section{Correlation and Linear Regression}


%choose soli deo gloria
\begin{center}
	\emph{-- Soli Deo Gloria --}
\end{center}


\newpage
\thispagestyle{empty}

\vspace*{\fill}
\begin{center}
	\emph{Soli Deo Gloria}
\end{center}
\vspace*{\fill}



\end{document}
